% RESEARCH_DRAFT_SHA256: f37f4b790150a6373c42171bfd1f3b716635df188d011c37de925bb0d6c108eb
% Layout and source-reference wrappers expanded for the web renderer.
\documentclass[11pt]{article}
\usepackage[margin=25mm]{geometry}
\usepackage{fontspec,unicode-math}
\setmainfont{Latin Modern Roman}
\setmathfont{Latin Modern Math}
\usepackage{amsmath,amssymb,mathtools,xurl,hyperref}
\hypersetup{colorlinks=true,urlcolor=blue,linkcolor=blue}
\setlength{\parindent}{0pt}
\setlength{\parskip}{5pt}
\setlength{\emergencystretch}{3em}
\newcommand{\sref}[2]{\hyperlink{source:#1:#2}{[S#2]}}
\newcommand{\cataloguescope}[1]{\paragraph{Source question.}#1}
\newcommand{\floor}[1]{\left\lfloor#1\right\rfloor}
\begin{document}

\section*{Unwinding series: a Hardy-space projection proof and the remaining boundary-rate question}
\noindent catalogue entry; catalogue entry. Research attempt dated 11 October 2026.
\paragraph{Result in scope.}This manuscript contains a source-specific attempt and complete proofs of the stated auxiliary calculations. It makes no new solution or novelty claim for the entire target.
\section{Original problem and preserved definition}
The following definition and source list are reproduced from the repository. Any earlier inspection dates and reported literature claims in this preserved text describe that source record, not additional verification by this attempt.
% BEGIN PRESERVED DEFINITION
\subsection{Definitions and mathematical statement}

\paragraph{Shared notation.}Standard number systems are denoted by $\mathbb N,\mathbb Z,\mathbb Q,\mathbb R,\mathbb C$; index ranges, computational models and specialized notation follow the cited problem source and the definitions below.


For a holomorphic function $F$ on the unit disk, write $F=B_0G_0$, extracting a Blaschke factor $B_0$ for its zeros. Iteratively factor $G_j-G_j(0)=B_{j+1}G_{j+1}$. The resulting unwinding expansion is formally $F=\sum_{j\geq0}G_j(0)B_0\cdots B_j$, with remainder obtained from the unfixed final $G_j$. Blaschke factors are inner factors associated with zeros in the disk.
\paragraph{Statement recorded in the supplied source.}
What can be rigorously established about convergence of the unwinding series, and can one obtain effective convergence rates for sufficiently regular functions $F$? \href{https://doi.org/10.4171/owr/2021/55}{[S2]}
\paragraph{Scope qualification.}
The regularity class, convergence norm, and zero-factor conventions must be specified in a theorem. They are deliberately not fixed by the supplied research question.
\paragraph{Source status.}
The supplied catalogue labels this record partially solved. Its source-status metadata was checked against the supplied snapshot on 2026-10-11. This import does not independently establish current open status or mathematical correctness. The catalogue literature assessment is dated 2026-08-22.
\subsection{Short English statement}
Prove convergence and quantitative approximation rates for the iterative zero-factor expansion.
\subsection{Sources}
\begin{itemize}
\item[S1] Ulam AI and the UnsolvedMath Contributors, UnsolvedMath: A Curated Collection of Open Mathematics Problems, catalogue entry. Catalogue curation and metadata: CC BY 4.0. The supplied source question is retained; definitions and the English summary are editorial additions.
\url{https://huggingface.co/datasets/ulamai/UnsolvedMath}.
\item[S2] Applied Harmonic Analysis and Data Science (2022), Oberwolfach Reports, DOI 10.4171/owr/2021/55, pp. 3007--3066
\url{https://doi.org/10.4171/owr/2021/55}.
\item[S3] Applied Harmonic Analysis and Data Science, Oberwolfach Reports 55/2021, primary report PDF
\url{https://ems.press/content/serial-article-files/46933}.
\textit{Underlying Oberwolfach Reports material retains its source-specific licence, including CC BY-SA 4.0 where applicable; see THIRD\_PARTY\_NOTICES.md.}
\end{itemize}
% END PRESERVED DEFINITION
\subsection{Verbatim supplied catalogue question}
What can be rigorously established about convergence of the unwinding series, and can one obtain effective convergence rates for sufficiently regular functions $F$?

\section{Source norm and convention choices}
The primary contribution, printed pages3020--3022, starts with
functions holomorphic near the closed unit disk, describes
Blaschke zero extraction, and asks about convergence and rates
without fixing the convergence norm. We give a complete
known-style $H^2$ projection argument and effective bounds with
explicit regularity constants. The 2023 primary abstract on
reproducing-kernel expansions already reports norm convergence
results in this direction; no novelty is claimed. Boundary
pointwise or uniform convergence is not inferred from an $H^2$
calculation.

\section{Finite-step identities and orthogonality}
Let $F$ be holomorphic in a neighbourhood of the closed unit
disk. If it is identically zero the expansion terminates
immediately. Otherwise extract all its finitely many zeros in
the open disk, including multiplicities, into a finite Blaschke
product $B_0$ and put $F=B_0G_0$. At each nonterminal step
set $c_j=G_j(0)$ and factor
\[
 G_j-c_j=B_{j+1}G_{j+1},
\]
where $B_{j+1}$ extracts all zeros in the disk. In particular
$B_{j+1}(0)=0$. If $G_j-c_j$ is identically zero, stop and
set all later remainders zero. Unimodular factor conventions
may change coefficients and factors but not these identities.
All finite-step quotients are holomorphic near the closed disk:
the divisions at the extracted zeros are removable.

Use $H^2=\{\sum_{k\ge0}a_kz^k:\sum_k|a_k|^2<\infty\}$
with the coefficient inner product. Multiplication by a finite
Blaschke product is an isometry. For polynomials this follows
from the boundary integral norm and $|B(e^{it})|=1$; for all
$H^2$ functions use polynomial density and extend the isometry
by completion. The extension agrees with analytic multiplication,
since point evaluation is bounded by
$|h(z)|\le\|h\|_{H^2}/\sqrt{1-|z|^2}$, directly from
Cauchy--Schwarz applied to the power series.

Write $\psi_j=B_0\cdots B_j$. Iterating the recursion gives
\[
 F=\sum_{j=0}^N c_j\psi_j+R_N,
 \qquad R_N=\psi_{N+1}G_{N+1}.
\]
For $j\le N$ and arbitrary $h\in H^2$, cancellation of the
common inner factor in the inner product gives
\[
 \langle\psi_j,\psi_{N+1}h\rangle
 =\langle1,(B_{j+1}\cdots B_{N+1})h\rangle=0,
\]
because the analytic function on the right vanishes at zero.
Complex-conjugate inner-product conventions yield the same zero.
Thus the partial sum is orthogonal to the closed subspace
$K_N=\psi_{N+1}H^2$. This subspace is closed because it is the
range of an isometry from a complete space. Since $R_N\in K_N$,
the identity proves that $R_N$ is exactly the orthogonal
projection of $F$ onto $K_N$, without an unproved completeness
claim about the functions $\psi_j$.

\section{A quantitative comparison with Taylor tails}
Each $B_j$, $j\ge1$, has a zero at zero, so
$K_N\subseteq z^{N+1}H^2$. If $F(z)=\sum_{k\ge0}a_kz^k$
and $T_NF=\sum_{k=0}^Na_kz^k$, every element of $K_N$ is
orthogonal to $T_NF$. Consequently
\[
 \|R_N\|_{H^2}^2=\langle F,R_N\rangle
 =\langle F-T_NF,R_N\rangle
 \le\|F-T_NF\|_{H^2}\|R_N\|_{H^2},
\]
where taking absolute values gives the same inequality under
either inner-product convention. Dividing when the remainder
is nonzero gives the explicit estimate
\[
 \|R_N\|_{H^2}\le
 \left(\sum_{k>N}|a_k|^2\right)^{1/2}.
\]
The right side tends to zero. This proves $H^2$ convergence
for the stated analytic-neighbourhood class, and in any more
general $H^2$ setting where the displayed inner factorizations
and recursions exist. Such existence is not proved here for
arbitrary boundary-singular functions.

Two effective rates follow. If
$\sum_{k\ge0}(1+k)^{2s}|a_k|^2\le C_s^2$ for $s>0$,
then $\|R_N\|_{H^2}\le C_s(N+2)^{-s}$ by bounding each
tail weight below by $(N+2)^{2s}$. If $F$ is holomorphic on
a neighbourhood of $|z|\le R$ with $R>1$ and
$\sup_{|z|=R}|F(z)|\le M$, Cauchy's coefficient formula gives
$|a_k|\le MR^{-k}$. Summing the geometric series yields
\[
 \|R_N\|_{H^2}\le
 \frac{M R^{-(N+1)}}{\sqrt{1-R^{-2}}}.
\]
Both the radius and the bound $M$ are needed for this effective
estimate; their existence alone is not an algorithm to compute
them from arbitrary data.

There is also a compact-interior estimate independent of such
extra regularity. Since $R_N=z^{N+1}h_N$ with
$\|h_N\|_{H^2}=\|R_N\|_{H^2}\le\|F\|_{H^2}$, for
$|z|\le r<1$,
\[
 |R_N(z)|\le\frac{r^{N+1}}{\sqrt{1-r^2}}\|F\|_{H^2}.
\]
This proves uniform convergence on every smaller closed disk.

\section{Why this route does not prove a boundary-uniform rate}
The implication from small $H^2$ norm to small supremum on the
unit circle is false. For any integer $L\ge1$, the polynomial
\[
 h_L(z)=L^{-1}\sum_{j=0}^{L-1}z^j
\]
has $H^2$ norm $L^{-1/2}$ but $h_L(1)=1$.
Multiplying by arbitrarily high powers of $z$ preserves both
facts and can enforce any desired order of vanishing at zero.
These polynomials are not claimed to be actual unwinding
remainders for one fixed function; they show that the norm
and zero-order estimates used above cannot alone imply the
boundary supremum conclusion. Additional control of the
remainders, their factors, or a stronger norm would be needed.

\section{Audit and status}
The audit checked the indexing of the remainder, termination at
a constant quotient, closedness of the inner-multiplier range,
and the zero at every subsequent step. In particular, mere
bounded residual energy was not used to infer strong convergence;
the orthogonal-projection identity supplies the missing step.
No root-finding algorithm, pointwise boundary theorem, numerical
experiment, or theorem prover is claimed. The complete scoped
norm result is known material rederived for this attempt; the
source's broader unspecified convergence/rate agenda is not
claimed resolved.

\section{Process and contributions}
Mathematics and model-assisted adversarial self-audit: OpenAI gpt-6-astra, ultra reasoning, as dispatched for this research batch; exact serving revision not exposed. Codex (GPT-6; exact serving revision not exposed) materially assisted with orchestration and format review. No human mathematical authorship, mathematical edits, or independent human proof checking is asserted.
The coordinating Codex (GPT-6; exact serving revision not exposed) also provided a separate model-assisted written-proof review of this manuscript. This is not independent human verification.
Contribution prompt A was applied to the full definition and supplied source row, then prompt B was applied adversarially to the derived results. The item-specific audit above records the scope checks and rejected extrapolations. Local assembly and structural checks do not establish mathematical correctness. The source catalogue status is preserved in TOP\_PROBLEM and is separate from this attempt status.
\paragraph{Integration review.} Codex (GPT-6; exact serving revision not exposed) reviewed the written arguments and expanded layout and source-reference wrappers. Repository metadata and the production web renderer were checked using Python 3.13 and JavaScript with MathJax 3.2.2. This model-assisted review is not independent human proof verification; no PDF compilation or formal proof checking is asserted.

\section{Completion Estimate}
\noindent\textbf{COMPLETION: 0\%}
The proved material is a rederivation of known or auxiliary facts rather than a substantiated new advance on the original target. The proof gives known-style H² and compact-interior convergence with explicit norm/radius hypotheses. It does not establish general boundary pointwise or uniform rates, effective zero-factor computation, or all unspecified regularity/norm regimes in the source.
\section{References}
The complete original Sources subsection above is retained, including catalogue attribution and licence notices. The additional list records primary sources actually consulted for this attempt on 11 October 2026; uncited source links reproduced above were not all rechecked.
\begin{itemize}
\item OWR55/2021, Applied Harmonic Analysis and Data Science; complete relevant primary passages read in cached report, with printed pages identified in the manuscript. \url{https://ems.press/content/serial-article-files/46933}.
\item Mashreghi and Verreault, Nonlinear expansions in reproducing kernel Hilbert spaces; primary abstract inspected for existing norm-convergence context, not used as an external proof dependency. \url{https://arxiv.org/abs/2310.01269}.
\end{itemize}
\section{Final status}
\textbf{Unresolved in this attempt.} The proof gives known-style H² and compact-interior convergence with explicit norm/radius hypotheses. It does not establish general boundary pointwise or uniform rates, effective zero-factor computation, or all unspecified regularity/norm regimes in the source.
\end{document}
